\section{Assumptions and Logbook}

\subsection{Assumptions and Dependencies}

Each is written so it could be shown false.
\begin{itemize}
\item \textbf{A1:} The watch's Wi-Fi radio can reach both the
custom board and the companion computer from the same room, with
no wired connection.
\item \textbf{A2:} Five volunteers performing five minutes of each
activity give enough labelled data to train a classifier that fits
the memory in C1.
\item \textbf{A3:} Motion features such as magnitude and cadence
are enough to separate sedentary, walking and running.
\item \textbf{A4:} The vendor wiki figures in C1--C4 match the
datasheet (not yet checked).
\item \textbf{D1:} A companion computer is available and running
for every link-degradation test.
\item \textbf{D2:} Each volunteer gives consent before any
recording.
\end{itemize}

\subsection{AI Use}

An AI assistant was used for the steps below, each with how the
output was checked.
\begin{itemize}
\item Explaining the required proposal structure and drafting
wording for each section. Checked: each section compared with the
instructor's skeleton; scope, needs and targets carried over from
the Project Idea Form.
\item Drafting the requirements table and the uncertainty
register. Checked: each requirement traced to a need and given a
method and acceptance criterion before any measurement.
\item Suggesting the 15\%, 10\,s and 30\,s targets and the
degraded-link and outage test conditions. Checked: accepted as
design choices with the reasons stated in Sections~\ref{sec:reqs}
and~\ref{sec:poc}; they are not published standards.
\item Typesetting in LaTeX. Checked: compiled in Overleaf and the
output read through.
\end{itemize}

\subsection{Logbook}

The decisions in this proposal are recorded in the logbook on
pages 1 to 11 (September 18 and 25, 2026) and in the catch-up
entry beginning on page 12 (October 15, 2026), which records the
decisions made between September 28 and October 15.
